\chapter{离散时间中的主动推断}

\begin{quote}
我不能创造的东西，我就无法理解。

--- Richard Feynman
\end{quote}

\section{引言}

到目前为止，我们一直在相对抽象的层面讨论主动推断的原理。本章转向具体例子，以及如何在实践情境中对这些例子加以规定。我们聚焦于离散时间中的类别变量模型。通过一系列复杂度逐步提升的例子，我们将说明感知处理、决策、信息搜寻、学习以及层级推断的模型。这些例子的选择目标，是尽可能简单地凸显主动推断方案中涌现出来的性质，包括可测量的生理表现与行为表现。

\section{感知处理}

我们首先考察感知处理，以及如何反演第 4 章引入的那类离散时间模型。稍后在本章中，我们会逐步构建一个完整的部分可观测马尔可夫决策过程（POMDP）。不过，我们先从 POMDP 的一个特殊情形开始：在这一情形下，我们可以忽略选择与行为，只关注隐藏马尔可夫模型（HMM）；它可用于顺序性的、类别性的感知推断（见图 7.1）。为此，我们先借助一个简单例子。想象你正在聆听一段短小乐曲的演奏。乐谱中写下来的音符序列可以看作隐藏状态（不可观测状态），而我们实际听到的音符序列则是结果（可观测结果）。如果演奏者是一位专业音乐家，那么隐藏状态与结果之间的对应关系会非常紧密。但如果演奏者是业余的，那么从本应弹出的音符到实际听到的音符之间，在似然映射上就会多出一层随机性。即便如此，只要我们对每个音符前后衔接另一个音符的概率具有先验信念，我们依然可能据此推断出原本应该听到的是哪个音符。

将聆听业余音乐家的这个例子形式化，方法如下。首先，我们要决定这位音乐家实际弹出她意图弹出的音符（结果）有多可靠。我们可以用 $A$ 矩阵来表达这一点，其中的元素表示在给定某个状态（列）时产生某个结果（行）的概率。在这个玩具例子中，我们将其设定为：

\begin{equation}
A=\frac{1}{10}
\begin{bmatrix}
7 & 1 & 1 & 1 \\
1 & 7 & 1 & 1 \\
1 & 1 & 7 & 1 \\
1 & 1 & 1 & 7
\end{bmatrix}
\label{eq:chapter7-A}
\end{equation}

这意味着有 $70\%$ 的时间里，这位音乐家会准确地弹出她打算弹的音符。接着我们规定 $B$ 矩阵中的转移概率，它刻画在给定当前状态（列）的情况下，下一个状态（行）的概率：

\begin{equation}
B=\frac{1}{100}
\begin{bmatrix}
1 & 1 & 1 & 97 \\
97 & 1 & 1 & 1 \\
1 & 97 & 1 & 1 \\
1 & 1 & 97 & 1
\end{bmatrix}
\label{eq:chapter7-B}
\end{equation}

这表示第一个音符之后有 $97\%$ 的概率接第二个，第二个之后有 $97\%$ 的概率接第三个，以此类推。如果我们知道这个序列总是以第一个音符开始，那么先验概率可设为：

\begin{equation}
D=
\begin{bmatrix}
1 \\
0 \\
0 \\
0
\end{bmatrix}
\label{eq:chapter7-D}
\end{equation}

式 \eqref{eq:chapter7-A}--\eqref{eq:chapter7-D} 共同完整规定了图 7.1 所示的 HMM 生成模型。换句话说，它们描述了我们关于“所听到的音乐如何由这位业余音乐家生成”的信念。利用式 4.12，并把这里的生成模型代入进去，我们就可以模拟一串结果所诱发的贝叶斯信念更新动力学。结果如图 7.2 所示。注意左上图中，随着时间推移积累到更多数据，置信度不断提高；但在第三个时间步出现了一个意外结果，因此这一趋势被打断。这个结果可以有两种解释。第一种是：在式 \eqref{eq:chapter7-B} 所表达的先验信念之下，被意图演奏的音符本身就是一个不寻常的音符。但由于该模型的 $B$ 矩阵中这类转移极其罕见，因此这种解释并不太可能。另一种、更可信的解释是：音乐家失误了，弹错了音。正如右上图第三列所示，我们模拟出来的听者最终采取的就是这种解释。不过，他仍然给“其实本来就是正确音符”保留了非零概率。能够报告这种不确定性，正是主动推断所提供的贝叶斯视角的一项关键特征。

这里展示的模型还可以在许多方向上变得更复杂，而也许最简单的一种方式就是对状态空间进行因子分解（Mirza et al. 2016）。例如，可以把音符分解为音高与力度（结果也做类似区分）。在视觉推断任务中，这种分解可以是“是什么”和“在哪里”，而这一区分在神经生物学中非常重要（Ungerleider and Haxby 1994）。在后续各节中，我们会利用这种因子分解，把那些受生物体自身影响的状态与那些不受其影响的状态区分开来。若想进一步了解这种不涉及行动的模型，以及可通过最小化自由能来反演它的神经元消息传递方案，可参见 Parr, Markovic et al. (2019)。

\section{作为推断的决策与规划}

上面的 HMM 说明了一种非常简单的、基于结果序列的类别推断形式。然而，这种模型所描绘的“固着不动”的生物相当乏味。自主生物显然不只是被动接受感觉数据的容器。相反，它们会主动改变环境，并与其感觉系统进行双向交换。这说明有必要把 HMM 扩展为 POMDP：我们不仅要推断环境如何变化，还要推断我们所选择的行动路径如何改变环境，以及应该选择哪一条行动路径。

图 7.3 展示了一个 POMDP 生成模型。这与第 4 章中介绍的模型相同，那里已经详细拆解了在这类模型中如何进行推断。请注意，它在结构上与图 7.1 中的 HMM 极其相似，只是多了一个额外变量 $\pi$，转移概率 $B$ 以它为条件。这意味着我们可以同时考虑关于状态动力学的不同假设。这些假设可以解释为生物体可在其间进行选择的计划。从这个角度看，策略评估等同于模型比较，而策略本身不过是对一串已观测到的（由自身生成的）感觉序列的解释变量。

图 7.3 中的模型与第 4 章介绍的版本有一个细微差别：它允许对状态（上标 $n$）和结果（上标 $m$）进行因子分解。当我们考虑视觉世界可被分解为物体的位置与物体的身份时，这种做法的效用就显而易见。若要表示“位置”和“身份”的每一种可能组合，不仅极其低效，而且会付出很高的复杂性代价；而通常情况下，身份不随位置变化，位置也不由身份决定。对于把时间与身份、位置分离开来，也可以提出类似论证（Friston and Buzsaki 2016）。在这个阶段引入因子分解的好处在于，我们可以把世界中那些受生物体控制的状态，与那些不受其控制的状态分离开来。前一类状态的转移概率会随策略不同而变化，而后一类状态则对策略保持不变。

有了这些预备知识之后，我们现在概述一个简单任务的例子（Friston, FitzGerald et al. 2017）。这个任务需要规划，并展示了使用 POMDP 的主动推断的一些关键方面。任务情境是一只处在 T 迷宫中的老鼠：一条臂上有厌恶性刺激，另一条臂上有吸引性刺激，而最后一条臂上则有一个线索，用来指示前两种刺激各自所在的位置。这样的设置意味着，老鼠的行为大致有两种方式。它可以直接走向两条可能藏有吸引性刺激的臂之一，但要承担遭遇厌恶性刺激的风险。或者，它也可以先去寻找信息性线索，然后再走向最有可能包含吸引性刺激的那条臂。

这个选择触及了心理学中的经典“探索---利用”困境，而在主动推断框架下，这一困境会被化解。其化解方式来自于：关于策略的先验信念要求最小化期望自由能。简要回顾一下（详见第 4 章），对于一个最小化变分自由能的生物体而言，最可能的策略，就是那些会导向最低期望自由能的策略。期望自由能具有如下形式：

\begin{equation}
G(\pi)=\mathbb{E}_{Q(\tilde{s}\mid \pi)}\!\left[H\!\left[P(\tilde{o}\mid \tilde{s})\right]\right]-H\!\left[Q(\tilde{o}\mid \pi)\right]-\mathbb{E}_{Q(\tilde{o}\mid \pi)}\!\left[\ln P(\tilde{o}\mid C)\right]
\label{eq:chapter7-G}
\end{equation}

这里第一、二项对应负的认知价值（即 $-I(\pi)$），最后一项对应务实价值。

把期望自由能分解为认知价值与务实价值，突出了两种驱动力：其一是朝向信息采集的（认知）驱动，其二是朝向实现先验信念的（务实）驱动（图 7.3 中的 $C$）。下一节我们会尝试更深入地建立对认知价值的直觉理解；暂时可以把它简单理解为，在某个特定策略下我们有望获得的信息量。务实价值的形式，实质上是把对所有策略取平均后的结果概率当成一种先验。这样一来，那些后果与该先验一致的策略就会变得更可能，因为它们对应更低的期望自由能。换一种更直观的说法：如果我们认为某类观察很可能会发生，那么我们就会采取行动去实现“自己将会遭遇这些观察”的信念。因此，结果的对数概率可以被看作与其他形式化框架中的效用函数等价，例如最优控制理论和强化学习中的效用函数。由于效用与信息价值都自然地作为期望自由能的两个组成部分涌现出来，我们就不再需要额外操心如何平衡探索与利用，因为二者都服务于同一个目标函数的优化。

要看这一点在 T 迷宫例子中如何展开，我们需要像前面的 HMM 那样，把生成模型形式化。图 7.4 至图 7.6 展示了构成 T 迷宫生成模型的似然与转移概率。我们会稍微详细地解释这一过程，因为这个最小例子提供了构建读者自己生成模型的基本积木。第一步，是决定有多少个结果模态来表示模型需要解释的（感觉）数据。这会告诉我们需要规定多少个 $A$ 矩阵。在这里，我们有两个模态：一个表示与老鼠在迷宫中所处位置有关的外感受数据（$A^1$），另一个是“是什么”的模态，可以理解为当老鼠找到吸引性（可食用）刺激时所经历的内感受数据（$A^2$）。这些模态中的各个水平，也就是每一种可能出现的观察结果，决定了每个 $A$ 矩阵的行数。下一步，要决定用于解释这些数据的隐藏状态因子的数量；这决定了我们需要多少个 $B$ 矩阵。这里我们考虑两个因子：老鼠在迷宫中的位置，以及情境（吸引性刺激在左边还是右边）。它们分别有四个与两个水平。然后，我们必须对每一种隐藏状态组合，规定各个结果出现的概率。图 7.4 给出了情境 1，图 7.5 给出了情境 2。对于第一个模态，$A^1$ 以概率 1 把每个位置映射到一个结果。线索位置会根据情境不同而与一个“向左线索”或“向右线索”相关联。第二个模态即内感受模态（$A^2$），它把起始位置与线索位置都关联到中性结果；当情境与老鼠所进入的迷宫臂一致时，它以 $98\%$ 的概率产生吸引性结果。严格来说，这些 $A$ 矩阵其实是张量对象，因为它们的元素由三个数字决定（结果、位置、情境），而矩阵只需要两个数字（行与列）即可确定。

随后我们需要规定转移概率。$B$ 矩阵描述在给定策略 $\pi$ 的情况下，从某一状态（列）转移到另一状态（行）的概率。它们规定了老鼠在迷宫中位置的转移（$B^1$），以及情境的转移（$B^2$）。图 7.6 给出了可控制的 $B^1$ 转移。每个矩阵对应一种不同的动作选择（以下标标记）。这些动作允许老鼠从任何位置移动到任何其他位置，惟独不能从迷宫的两条臂离开，因为那两条臂是吸收状态。这意味着一旦进入其中，无论老鼠再选什么动作，它都必须停留在那里。相比之下，老鼠无法控制情境，也就是说，它无法控制自己究竟处于图 7.4 所示的情境 1，还是图 7.5 所示的情境 2。情境在时间上保持不变，因此可以表示为单位矩阵：

\begin{equation}
B^{\pi 2}=
\begin{bmatrix}
1 & 0 \\
0 & 1
\end{bmatrix}
\label{eq:chapter7-context}
\end{equation}

这里每一列（和每一行）都对应于图 7.4 或图 7.5 中的一个状态索引。这意味着无论起始于哪个情境，它都会随时间保持不变，也即只发生对自身的转移。而且，不管选择什么策略，这一点都成立。向量 $C^1$ 给出了该模态下各结果的先验偏好：总体上是均匀偏好，但对起始位置有轻微厌恶（$-1$）。向量 $C^2$ 则表达了对吸引性刺激的偏好（$+6$）以及对厌恶性刺激的排斥（$-6$）。对二者都没有遇到则视为中性（$0$）。

\begin{equation}
\begin{aligned}
C^1 &= \sigma\!\left([-1,0,0,0,0]^T\right) \\
C^2 &= \sigma\!\left([0,6,-6]^T\right)
\end{aligned}
\label{eq:chapter7-C}
\end{equation}

这些向量中元素的顺序，对应于各自 $A$ 矩阵中行的顺序。Softmax 函数 $\sigma$ 允许我们用正值和负值来表述偏好（对应未归一化的对数概率），再把它们转换为概率。这既保留了对数概率之间的差异（或相对概率），又保证了归一化。从实践上看，这意味着在老鼠的生成模型中，吸引性刺激被认为比中性刺激的概率高出 $e^6$（约 $400$）倍。这是一种非常强烈的偏好，意味着老鼠相信自己的行动更有可能导向吸引性结果。正是这种对行动推断的约束，对接下来出现的行为至关重要。最后，$D$ 向量规定初始状态的先验概率：

\begin{equation}
\begin{aligned}
D^1 &= [1,0,0,0]^T \\
D^2 &= \frac{1}{2}[1,1]^T
\end{aligned}
\label{eq:chapter7-D2}
\end{equation}

这些向量中元素的顺序与 $B$ 矩阵中的顺序一致。$D^1$ 表示对“从迷宫中心开始”的高度确信；$D^2$ 则表示在起始时，两种情境（图 7.4 或图 7.5）被认为等可能。

图 7.7 展示了当我们对图 7.4 至图 7.6 的生成模型进行反演时会发生什么。上排描述的是如果我们直接观察老鼠行为所能看到的情形。老鼠先从中心出发，然后前往信息性线索所在位置。这是因为该位置具有很高的认知价值，也就是说，在这个位置获得的观察有潜力消除关于情境的不确定性。看到指向左侧情境（情境 1）的线索后，老鼠便选择迷宫左臂，并在那里找到奖励性刺激。这一步是由该位置的高务实价值驱动的。下方图像展示了这一简单试次中发生的信念更新。与图 7.2 一样，这里也以一种理想化老鼠的可观测形式来表示这些更新，也就是说，如果我们测量神经活动（如发放率和局部场电位，LFP），理论上会看到这样的形式。请注意，在第二个时间步，也就是老鼠到达信息性线索位置时，信念发生了快速变化，并伴随出现相应的 LFP（虚线）。

\section{信息搜寻}

第 7.2 节中的模拟展示了一个简单的探索---利用权衡例子：先通过搜寻信息来消除不确定性，然后再利用已推断出的内容去满足先验偏好。本节将更细致地剖析认知价值这一概念。正如式 \eqref{eq:chapter7-G} 所示，它由两项组成：

\begin{equation}
\begin{aligned}
I(\pi) &= H[Q(\tilde{o}\mid \pi)]-\mathbb{E}_{Q(\tilde{s}\mid \pi)}[H[P(\tilde{o}\mid \tilde{s})]] \\
&= D_{\mathrm{KL}}\!\left[P(\tilde{o}\mid \tilde{s})Q(\tilde{s}\mid \pi)\,\middle\|\,Q(\tilde{o}\mid \pi)Q(\tilde{s}\mid \pi)\right] \\
&= \mathbb{E}_{Q(\tilde{o}\mid \pi)}\!\left[D_{\mathrm{KL}}\!\left[Q(\tilde{s}\mid \pi,\tilde{o})\,\middle\|\,Q(\tilde{s}\mid \pi)\right]\right]
\end{aligned}
\label{eq:chapter7-I}
\end{equation}

这三种写法分别突出了不同视角：第一行是后验预测熵减去期望模糊性；第二行把它写成互信息；第三行则把它写成信息增益，也就是显著性或贝叶斯惊奇。为了更直观地理解这些量，我们不妨借用视觉范式来说明：不同的眼跳（$\pi$）会导致不同注视位置（$s$）之间的转移。除了注视位置之外，隐藏状态还包括各位置上刺激的身份。刺激身份与注视位置共同生成视觉和本体感觉上的后果（$o$）。

在这个框架下，后验预测熵可以理解为“如果我执行这次眼跳，我会看到什么”这一问题所对应的不确定性程度。从科学家的角度看，它衡量的是：如果我们执行某个实验，我们对将会获得的数据究竟有多不确定。沿着这个思路，选择那些后验预测熵最大的眼跳（或实验）是合理的，因为这类行为拥有最大的潜力去消除不确定性。如果我们已经高度确信某个实验会得到什么结果，那么去做它就几乎没有收益。

然而，预测熵只告诉我们总共有多少不确定性，并没有告诉我们其中有多少是真正可以被消除的。对于一串随机生成的数字，我们总会对下一个数感到不确定，但即便一直盯着它们看，我们也永远无法消除关于其生成过程的不确定性。这正是期望模糊性发挥作用的地方。它衡量的是观察结果与状态彼此独立的程度。如果状态总是生成同一个观察，这个量就为零；而如果像随机数生成器那样，状态与结果之间毫无关联，它就达到最大。在视觉领域中，这意味着最好的眼跳目标将是一个照明良好的刺激，因为对“如果我去看这个刺激，会看到什么”几乎没有模糊性。综合来看，最好的眼跳，也即最好的感知实验，是那些既存在大量不确定性有待消除（后验预测熵高），又确实可以被消除（模糊性低）的情形。有趣的是，这与统计学中根据信息增益来评估实验设计的表达式具有完全相同的形式（Lindley 1956）。

图 7.9 展示了在一个眼跳范式中（Parr and Friston 2017b），当我们操纵模糊性和后验预测熵时会发生什么。图中有四个刺激（方块），每个方块的颜色都会随时间变化。其上叠加了一条模拟的眼动轨迹，仿佛我们正在测量实验参与者注视的位置。关键在于，我们把对结果的先验信念设为均匀，也就是让务实价值不存在，从而排除了任何基于偏好的选择。这样一来，每一次眼跳都是为了最大化认知价值而被选中。当生成模型把四个刺激都视为等价时（左图），四个位置会以大致相同的频率被采样。然而，我们可以调节各刺激所对应的不确定性（见框 7.1）。如果我们把某个刺激设得更模糊，即通过增大对应 $A$ 矩阵非对角元素的值，这个方块就会被忽略（中图）。这就是著名的“路灯效应”（streetlight effect）的一个例子（Demirdjian et al. 2005）。这个名称来自这样一个比喻：有人深夜丢了钥匙，最先找的地方往往是路灯下面，并不是因为钥匙最可能在那里，而是因为那里最容易获得高质量、低歧义、能够消除不确定性的信息。模拟表明，那个模糊的（例如光线昏暗的）方块会被忽略，从而在计算机模拟中复现了路灯效应。

相对地，图 7.9 右图展示了当我们让左下方方块的转移变得更难预测时会发生什么。我们会很快在该位置积累不确定性，因此在不改变模糊性的前提下，它拥有很高的后验预测熵。正如图中所见，这会导致眼睛更频繁地注视该位置，因为那里总有新的不确定性等待消除。直观来说，如果我知道某件东西的动力学非常可预测，那么我不必经常去看它，也可以对它的状态保持信心。相反，如果某件东西可能在我看别处的时候发生了变化，那么值得回头去检查。这些模拟的目的，是帮助我们建立对认知价值两部分的直觉理解，并看清期望自由能最小化如何确保我们主动选择感觉数据，以便更好地了解世界。

\subsection*{框 7.1：不确定性与精度}

图 7.7 中的例子诉诸了“精度”这一概念，而这是本书中的一个重要思想。精度是方差的倒数，用于刻画我们对某个概率分布的置信程度。它与一个分布的负熵（negentropy）密切相关：

\begin{equation}
-H[P(s)] = \mathbb{E}_{P(s)}[\ln P(s)]
\end{equation}

要给一个分布加上“可以变得更精确或更不精确”的参数化方式，最简单的方法之一就是采用带有逆温参数 $\omega$ 的 Gibbs 形式。其表达式为：

\begin{equation}
P(s\mid \omega)=\mathrm{Cat}(\sigma(\omega \ln D))
\end{equation}

注意，精度是乘在对数先验上的，因此它可以被理解为一种增益控制装置，即它会放大神经信号，而不是简单地把信号相加。图 7.8 中的图像展示了：对同一个 $D$，当我们改变 $\omega$ 时，概率分布会如何变化（每一列代表某个备选状态的概率）。可以看到，随着精度升高，置信度也在升高。

这种参数化方式可以应用到 POMDP 中用到的任何分布上。此外，我们还可以对精度本身定义先验，并像推断其他潜变量一样去推断它，即通过自由能最小化。若假定先验服从 Gamma 分布（从而排除精度取负值的可能），则可得到如下更新式（推导见附录 B）：

\begin{equation}
\begin{aligned}
P(\omega) &= \Gamma(1,\beta^\omega) \\
Q(\omega) &= \Gamma(1,\beta^\omega) \\
\Rightarrow \beta^\omega &= (D\beta^\omega - s)\cdot \ln D + \beta^\omega - \beta^\omega
\end{aligned}
\end{equation}

越来越多的研究认为，这些精度参数在生物学上的底层载体，可能正是那些调节神经反应增益的神经调质系统。第 5 章讨论了将这些参数与特定神经化学物质相联系的证据。

\section{学习与新颖性}

第 7.2 至 7.4 节已经给出了大多数主动推断实际应用所需要的一切。不过，在前面的讨论中，我们一直假定生成模型本身是已知的，并且不会因学习而发生改变。而在一些实际应用中，我们可能希望考虑生成模型的某些部分（例如 $A$ 矩阵或 $B$ 矩阵）如何在实验过程中被学习；更广泛地说，还希望考虑如何在已有数据的基础上优化生成模型的结构本身（Friston, FitzGerald et al. 2016）。一旦如此，主动推断就扩展为主动学习，而描述“关于状态的信息增益”的显著性（式 \eqref{eq:chapter7-context} 之后的讨论，或更准确地说式 \eqref{eq:chapter7-I} 中的量）就会得到一个补充：新颖性。新颖性关心的，是例如式 \eqref{eq:chapter7-A} 中 $A$ 矩阵元素、式 \eqref{eq:chapter7-B} 中 $B$ 矩阵元素，或生成模型其他参数上的不确定性如何被消除。此时，这些信念不再像之前假定的那样固定不变，而是可以随时间变化（Schwartenbeck et al. 2019）。为了达到这一点，我们首先必须像图 7.10 那样扩展生成模型，把关于这些模型参数的信念也包含进来。

从概念上说，把关于参数的信念纳入生成模型，就意味着可以把学习也看作另一种贝叶斯推断：也就是从关于模型参数的先验信念，过渡到后验信念。这凸显了感知与学习在根本上的相似性：就像感知可以描述为通过反演生成模型来根据观察推断隐藏状态一样，学习也可以描述为反演一个包含参数信念的生成模型，只不过这种反演通常运行在更慢的时间尺度上。

\subsection*{框 7.2：共轭先验}

在搭建图 7.10 这种形式的生成模型时，审慎选择适合先验信念的分布极其重要。通常，这应该是与似然相对应的共轭先验分布。所谓共轭先验，是指在进行贝叶斯推断时，后验信念会保持与先验相同的分布类型。举例来说，根据贝叶斯公式：

\begin{equation}
P(D\mid s)\propto P(D)\,P(s\mid D)
\end{equation}

如果 $P(s\mid D)$ 是类别分布，那么只要我们为 $P(D)$ 选择与类别分布共轭的 Dirichlet 分布，就可以保证 $P(D\mid s)$ 也仍是 Dirichlet 分布。形式化地写：

\begin{equation}
\left.
\begin{aligned}
P(D) &= \mathrm{Dir}(d) \\
P(s\mid D) &= \mathrm{Cat}(D)
\end{aligned}
\right\}
\Rightarrow P(D\mid s)=\mathrm{Dir}(d)
\end{equation}

要选择合适类型的先验，最简单的办法就是查找与当前似然分布形式相对应的共轭先验。对于本章使用的类别分布，适合作为参数信念的就是 Dirichlet 分布（见框 7.2）。引入这些额外的先验信念之后，我们现在就可以优化关于生成模型结构的后验信念。这意味着，要像第 4 章中处理状态那样，把这些量纳入自由能之中，并寻找自由能的极小值：

\begin{equation}
\theta=(A,B,C,D,E)
\label{eq:chapter7-theta}
\end{equation}

\begin{equation}
F=\mathbb{E}_{Q(\pi,\theta)}[F(\pi,\theta)]+D_{\mathrm{KL}}[Q(\theta)\|P(\theta)]+D_{\mathrm{KL}}[Q(\pi)\|P(\pi)]
\label{eq:chapter7-freeenergy-learning}
\end{equation}

Dirichlet 分布由计数（或伪计数）参数化，这些参数表示某个类别变量出现过多少次；对于先验而言，则表示“仿佛它曾被观测到这么多次”。关于这些参数的更新规则推导，可见附录 B。这里我们先概括其更新规则与 Dirichlet 分布的几个关键性质，重点关注与 $A$ 的先验和后验相联系的浓度参数 $a$ 与 $\bar{a}$：

\begin{equation}
\begin{aligned}
\bar{a} &= a + \sum_\tau s_\tau \otimes o_\tau \\
\mathbb{E}_Q[A_{ij}] &= \bar{A}_{ij} \approx \frac{a_{ij}}{a_{0j}} \\
\mathbb{E}_Q[\ln A_{ij}] &= \ln \bar{A}_{ij} = \psi(a_{ij})-\psi(a_{0j}), \qquad a_{0j}=\sum_k a_{kj}
\end{aligned}
\label{eq:chapter7-dirichlet}
\end{equation}

第一行表达的是：在经历了一系列观察之后，并结合对引起这些观察之状态的信念，浓度参数如何从先验更新为后验。圆圈中的叉号表示 Kronecker 张量积（当作用于两个向量时，就是外积）；在这里，它生成一个矩阵，其每个元素都等于 $s_\tau$ 和 $o_\tau$ 中某一对元素的乘积。这个更新规则可以被简单理解为一种活动依赖型可塑性：当某个结果被观测到，同时我们在后验上相信某个特定状态导致了这一结果，那么表示两者关系的矩阵元素就会增加。第二行强调了把 Dirichlet 浓度参数理解为“计数”的意义：对于给定状态（某一列），$a$ 中每个元素都表示相应结果出现过多少次。把它除以这一列元素之和（即总观测次数或总伪观测次数），就得到在该状态下各结果的概率。

要理解这种（伪）计数方法为什么直观，可以回到本章开头的业余音乐家例子。如果我们统计：当她本来想弹第一个音符时，实际弹出第一个音符多少次（第一行第一列）；想弹第二个音符时，实际弹出第二个音符多少次（第二行第二列）；依此类推；然后把这些计数除以她每个目标音符总共尝试过多少次，那么最终就会收敛到式 \eqref{eq:chapter7-A} 所示 $A$ 矩阵的正确数值，也就是她有 $70\%$ 的概率会正确命中自己想弹的音符。计数法还带来另一个重要后果，我们稍后还会回到这一点：在一次观察发生之前，已经积累了多少计数或伪计数，会决定我们在看到这个观察时多大程度上更新自己的信念。设想我们掷一枚硬币五次，结果连续五次都是正面，这可能会使我们倾向于认为这是一枚不公平的硬币；但如果在此之前已经掷过 100 次，并且其中 50 次正面、50 次反面，那么最后这连续五次正面，对我们关于“这枚硬币是否公平”的信念影响就会很小。式 \eqref{eq:chapter7-dirichlet} 的第三行给出了与 Dirichlet 分布有关的一个有用恒等式：随机变量对数的期望，可表示为两个 digamma 函数（Gamma 函数导数）之差。

这种以推断来理解学习的方式，突显了主动推断与大多数计算神经科学和机器学习方法之间的重要差异。后者往往引入各种学习规则，例如 Hebb 规则或误差反向传播，并把它们视作生物学上更真实或计算上更高效的机制；而在主动推断中，支配学习的更新规则是从统计学考虑中推导出来的，但它们最终却与生物学上关于活动依赖型可塑性的规则惊人地相似（见上面对式 \eqref{eq:chapter7-dirichlet} 第一行的讨论）。这正体现了规范性方法的一大吸引力：它从第一原理出发，不仅解释我们已经知道的大脑与行为事实，也解释那些我们之前并不知道的现象。

主动推断与大多数机器学习方法的另一个差别是：学习在这里天然被描述为一个主动过程，生物体会自主选择最有利于改进其生成模型的数据。当我们把关于参数的信念纳入模型时，这一点就变得尤其明显，因为期望自由能会多出一项：

\begin{equation}
\begin{aligned}
G(\pi) &= D_{\mathrm{KL}}[Q(\tilde{o}\mid \pi)\|P(\tilde{o}\mid C)] + \mathbb{E}_{Q(\tilde{s}\mid \pi)}[H[P(\tilde{o}\mid \tilde{s})]] \\
&\quad + \mathbb{E}_{Q(\tilde{o},\tilde{s},\theta\mid \pi)}[\ln Q(\theta)-\ln P(\theta\mid \tilde{o},\tilde{s})] \\
&= -\mathbb{E}_{Q(\tilde{o}\mid \pi)}\!\left[D_{\mathrm{KL}}[Q(\tilde{s}\mid \pi,\tilde{o})\|Q(\tilde{s}\mid \pi)]\right] \\
&\quad - \mathbb{E}_{Q(\tilde{o},\tilde{s}\mid \pi)}\!\left[D_{\mathrm{KL}}[Q(\theta\mid \tilde{o},\tilde{s})\|Q(\theta)]\right] \\
&\quad - \mathbb{E}_{Q(\tilde{o}\mid \pi)}[\ln P(\tilde{o}\mid C)]
\end{aligned}
\label{eq:chapter7-novelty}
\end{equation}

前两节已经有了显著性与务实价值项，但这里新增的是新颖性项。最后一个等式的写法凸显了显著性与新颖性之间的关系。简而言之，显著性之于推断，正如新颖性之于学习。二者都表达了这样一个思想：一旦执行了某个感知实验（也即策略中的一个动作），我们的信念会发生多大变化。正如科学实验一样，数据收集之后信念变化越大，这个实验就越“好”。

回到掷硬币并累积计数的比喻，它还能告诉我们另一件有用的事。假如我们有两枚硬币，而且可以任选其一来掷，那么如果想引发最大的信念变化，我们应该去掷那枚之前只掷过 5 次的硬币，而不是那枚已经掷过 100 次的硬币。前者具有更大的新颖性，因为我们对它还不熟悉。类似地，如果我们对某些变量的先验信念非常确信，就好像已经观测过它们很多次一样，那么去探查这些变量的策略所对应的新颖性就会比探查那些我们知之甚少的变量要低。

为了说明这一点在实践中如何运作，设想一个高度近视的生物站在铺有瓷砖的地板上。这个生物只能看到自己脚下那块砖的颜色，而且每次只能移动一块砖的距离。对于一个足够大、包含许多砖块的场景而言，如果把每块砖的颜色都表示为不同的隐藏状态，在计算上会非常昂贵。不过有一种更简单的模型可用：如果我们只把位置与隐藏状态相关联，把颜色只与结果相关联，那么就可以把“如果我走到那边，会看到什么”这样的信念高效地表示在一个由位置生成彩色砖块的 $A$ 矩阵中。通过累积 Dirichlet 参数（式 \eqref{eq:chapter7-dirichlet}），这个生物可以基于观察去优化这些信念。我们可以把这理解为一种突触记忆，而不是通过持续的神经活动来维持“某块砖是什么颜色”的信念。对于这种生成模型而言，所有不确定性都位于似然分布的参数中；于是，在没有任何偏好的情况下（也就是式 \eqref{eq:chapter7-novelty} 中新颖性主导策略选择时），看看会发生什么就很有意思。图 7.11 展示了一个由 64 块黑白砖组成的简单环境的模拟。每访问一块砖，关于该位置出现黑色还是白色的似然信念都会通过累积 Dirichlet 参数而得到更新。由于较大的 Dirichlet 参数会阻止信念发生大幅更新，期望自由能最小化所驱动的新颖性消解，会使我们的模拟生物避免回到任何之前已经访问过的位置。

同样的原理还可以应用到许多其他范式中。例如，如果我们把该生物的移动路径重新解释为一条眼跳扫描路径，那么这个框架就可以用于主动视觉采样。在主动视觉领域，这一框架已经被用来模拟由目标划消任务所诱发的视觉搜索行为（Parr and Friston 2017a）。随后，相关研究还展示了在这种情境下积累 Dirichlet 参数所需的短时程可塑性的证据（Parr, Mirza et al. 2019）。

正如我们能够把关于推断的思想扩展到学习一样，我们还可以至少再向前迈一步，去思考结构学习：它不只是优化模型中的参数，而是在参数多少不同的模型之间进行选择。框 7.3 给出了一种实现方式：它允许对不同假设模型进行高效的事后比较。这种思路已被用作睡眠的隐喻（Friston, Lin et al. 2017），以及静息态自发活动的隐喻（Pezzulo, Zorzi, and Corbetta 2020）：即便没有收集新的数据，模型的结构仍然可以被精炼和简化。

\subsection*{框 7.3：结构学习与模型约简}

第 7.5 节前半部分讨论的是一种重要但有限的（参数性）学习形式。进一步的复杂化，也就是学习世界结构本身，超越了对模型参数的优化，而转向这样一个问题：我们是否应该扩展模型结构，或者修剪模型结构？这可以被表述为模型比较问题（Friston, Lin et al. 2017）。换句话说，如果我消除似然矩阵中的某些元素，自由能会升高还是降低？通过比较含有这些元素和不含这些元素的模型，我们就可以回答这个问题。然而，显式地反演多个模型可能代价很高。幸运的是，有一种高效方法可以做到这一点，这就是贝叶斯模型约简（Friston, Litvak et al. 2016; Friston, Parr, and Zeidman 2018）；它只需要反演一个完整模型。在一般情形下，完整模型与采用另一组先验（记作 $\tilde{\ }$）的模型之间的比较，可通过下列公式实现：

\begin{equation}
\Delta F = F[\tilde{P}(\theta)]-F[P(\theta)] = \ln \mathbb{E}_{Q(\theta)}\!\left[\frac{\tilde{P}(\theta)}{P(\theta)}\right]
\end{equation}

\begin{equation}
\tilde{Q}(\theta)\propto \exp\!\left(\ln Q(\theta)+\ln \tilde{P}(\theta)-\ln P(\theta)+\Delta F\right)
\end{equation}

对第 7.5 节使用的 Dirichlet 先验而言，上式可写成如下形式（其中 $B$ 是多元 Beta 函数）：

\begin{equation}
\Delta F = \ln B(\tilde{d})-\ln B(d)+\ln B(a)-\ln B(\tilde{a})
\end{equation}

\begin{equation}
\tilde{a}=a+\tilde{d}-d
\end{equation}

这种模型约简形式，可能有助于理解离线状态下的模型优化，例如睡眠中可能发生的那种优化。第 8 章中，当我们讨论同时具有离散与连续成分的层级模型优化时，还会简要回到贝叶斯模型约简。

\section{层级推断或深度推断}

上一节中，我们看到了一种对原始生成模型进行层级扩展的方法，即为生成模型的参数定义先验。图 7.12 展示了第二种层级形式，它关涉时间尺度的嵌套。这个用于层级推断或深度推断的生成模型，可以被看作图 7.3 所示浅层模型的层级扩展：在较低层级上，它包含一系列与图 7.3 相同的 POMDP 模型（其中一个例子由虚线框标出），而这些低层模型又被更高一层的 POMDP 所情境化。

重要的是，这个生成模型包含了在不同时间尺度上演化的变量：高层级变化较慢，低层级变化较快（Friston 2008; Friston, Rosch et al. 2017; Pezzulo, Rigoli, and Friston 2018）。这一点可以从如下事实中看出：第 1 层的 POMDP 模型各自演化 3 个时间步，但这些短轨迹中的每一条，都依赖于高层级（第 2 层）上的一个单一状态，而这个高层状态在整个低层轨迹期间保持不变。换句话说，从高层级的视角看，每经过 1 个时间步，低层级实际上已经运行了多个（这里是 3 个）时间步。

为了对这种时间尺度分离，也就是支撑深度时间推断的机制，形成一些直觉，不妨考虑日常生活中的一个简单层级例子：阅读。我们对构成句子的词进行推断；句子又组合成段落；段落构成页面；页面构成书籍；书籍构成图书馆，以此类推。如果我们把图 7.12 低层级中的每个状态想象成一个词，那么高层级中的每个状态就可以看作一个句子。关键在于，句子的持续时间跨越了序列中任何单个词的持续时间。

阅读的例子在图 7.13 中有更详细的说明，它基于 Friston, Rosch et al. (2017) 的例子，更多细节可参见该文。该模型的结构与图 7.12 相同，表示一个非常简单语言中的句子（高层）与词（低层）。这个语言包含三个可能的词：\emph{flee}、\emph{feed}、\emph{wait}，它们可以组成六种可能的四词句子。如果句子是“flee, wait, feed, wait”，那么高层会对第一个低层 POMDP 预测词 \emph{flee}，对第二个低层 POMDP 预测词 \emph{wait}，依此类推。在低层，我们从一个基于高层的经验性先验 $D$ 出发，它告诉我们哪些词更可能出现。例如，如果我们在图 7.13 上方面板所示句子上采用均匀分布，那么会发现第一个词在其中三分之二的句子里是 \emph{wait}，其余三分之一则是 \emph{flee}。这意味着，在第一个低层 POMDP 的第一个时间步中，我们的 $D$ 向量就应当为这些词赋予相应概率。

低层的词随后会生成观察，也就是基于当前注视于词的哪一部分所形成的视觉输入。与图 7.9 的例子类似，这个 POMDP 允许选择不同的中央凹注视目标，从而为各个假设词积累支持或反对的证据。这同样诉诸前文已经介绍过的期望自由能最小化过程，因此我们这里不再详细展开具体的注视动作，只指出：随着低层级每一个时间步的推进，对当前词的置信度都会提高。在图 7.13 所示的序列中，我们看到在前几个快时间尺度 $\tau^{(1)}$ 的时间步里，低层为 \emph{flee} 这个词积累了证据。这一推断随后被向上传播到高层，于是为第一句和第四句提供支持，因为它们都以这个词开头。在后续时间步中，低层积累到的证据与这两个句子都一致。到了第四个慢时间尺度 $\tau^{(2)}$ 的步骤，在第一句之下我们会预测 \emph{wait}，在第二句之下则会预测 \emph{flee}。当在快时间尺度上推断出的是 \emph{wait} 时，慢时间尺度上的第一句便被推断出来。到最后一步，模拟成功选中了正确句子，并获得正确反馈作为奖励。由此产生的信念更新可以在图 7.12 下部的 LFP 图中看到。

这类深度时间模型已被用于模拟阅读（Friston, Rosch et al. 2017）、延迟期工作记忆任务（Parr and Friston 2017c），以及为视觉推断计算经验性先验（Parr, Benrimoh et al. 2018）。此外，它们还被用于关于动机与控制的理论解释（Pezzulo, Rigoli, and Friston 2018）。原则上，这类模型可以扩展到任意多的层级，从而刻画一个具有深度结构、并且在许多不同时间尺度上展开动力学的世界。

我们还可以在图 7.12 与图 7.13 那类层级模型，以及图 7.10 那类学习模型之间，看出一个有趣的平行关系。学习模型可以被视为层级生成模型，它强调的是较快的推断动力学（对状态信念的更新）与较慢的学习动力学（对参数信念的更新）之间的时间尺度分离。图 7.10 与图 7.12 所示的模型还可以按照任意复杂度进行组合；在这样的组合中，不同层级变量之间的关系本身也可以被学习。这使得我们能够设计出越来越复杂的生成模型，以应对系统层面的认知与神经生物学问题。

\section{总结}

本章展示了如何构造离散时间生成模型，以处理一系列认知与神经生物学问题，例如感知推断、决策与规划、探索与利用的平衡、参数学习与结构学习，以及新颖性搜寻。这当然远不能穷尽离散模型在主动推断中的全部应用，但它足以说明这类建模的关键原则。前面介绍的模型可以进行层级组合，可以加入参数上的先验，也可以为策略或偏好引入情境敏感的先验。更重要的是，无论是简单生成模型还是更复杂的生成模型，其推断都始终可以通过自由能最小化来推进，这恰好说明了这种方法的普适性。主动推断的不同方面会在不同生成模型下显现出来，例如当我们为模型参数引入先验时，新颖性搜寻就会浮现出来；这为通过设计合适的生成模型来探索一个开放无边的认知与生物学问题集合，打开了可能性。
